\section{Introduction}

\noindent Cellular networks are critical for modern global connectivity, with 5G and 4G forming the backbone of modern mobile communication systems. Technical specifications and architectures are standardized by the 3rd Generation Partnership Project (3GPP), an association of seven Organizational Partners, including ATIS in the U.S., ETSI in Europe, and CCSA in China. Thousands of specifications are published by 3GPP to ensure global interoperability across mobile networks and devices.
Consequently, mobile equipment vendors and operators are required to ensure compliance with 3GPP specifications. However, this standardization introduces systemic risks: 
\textit{Design flaws in 3GPP specifications can propagate into all compliant networks and devices, creating widespread threats with global negative impact}. 
Moreover, the increasing complexity, frequent revisions, and evolving new generations of mobile networks continue to expand attack surfaces. 

To mitigate those risks, researchers spend increasing effort to scrutinize cellular network specifications and services documented in the 3GPP specifications to identify design issues. However, the traditional manual analysis by experts is slow, error-prone, and infeasible at scale. 
State-of-the-art automated approaches primarily rely on rule-based approaches (e.g., dependency parsing \cite{chen2021bookworm,al2024hermes}), and keyword-driven scanning (e.g., hazard indicators~\cite{chen2021bookworm}), failing to capture the deeper and implicit knowledge in 3GPP specifications.  
As a result, vulnerabilities in critical mobile network procedures continue to be exploited, enabling attacks such as downgrade attacks~\cite{shaik2015practical, chen2021bookworm, 6550445, kambourakis2011attacks,
      lee2009detection, leong2014unveiling, hussain20195greasoner,
      hussain2018lteinspector, kim2019touching, 8894379}. 
These methods suffer from three fundamental limitations: (1) they lack deep insights into the logic, processes, and context of specification texts, particularly when dealing with cross-sentence/clause dependencies; (2) their security requirements are manually developed for specific procedures, which severely limits scalability across the thousands of 3GPP specifications; (3) their analysis pipelines require costly manual updates to remain synchronized with evolving 3GPP releases.  
These limitations highlight the urgent need for an intelligent, automated system capable of understanding 3GPP specifications and systematically assessing them.  

In this study, we develop \lm (\ul{Cell}ular Network \ul{Sec}urity \ul{Inspector}), an automated framework capable of understanding 3GPP specifications and systematically assessing them via foundational security properties. \lm models mobile networks in structured SCA (state-condition–action) representations from 3GPP specifications, and examines them with foundational security properties. By integrating specification analysis with automated test case generation, \lm delivers a scalable, adaptive, and future-proof solution for 3GPP specification security assessment. In this study, \lm analyzes the critical NAS and RRC layers for current operational 5G and 4G mobile networks as well as the selected sections of TS 23.501 and TS 24.229. 
\lm not only discovers 36 previously reported vulnerabilities but also reports 7 new vulnerabilities.
We summarize our contributions as follows. 

\noindent $\diamond$\textbf{New Framework:} We developed \lm, a novel approach for automated security assessment of 3GPP specifications, advancing toward fully automated and scalable specification analysis for safeguarding not only operational 5G/4G cellular networks but also Next-G mobile networks.  

\noindent $\diamond$\textbf{New Findings:} Applying \lm to 5G and 4G NAS and RRC specifications, as well as the selected sections of TS 23.501 and TS 24.229, we discovered 43 vulnerabilities, 7 of which were previously unreported. They are shown to have serious security consequences, demonstrating the urgent need for automated 3GPP specification analysis and the effectiveness of \lm in identifying known and previously undetected vulnerabilities. 
